% Part: first-order-logic
% Chapter: completeness
% Section: downward-ls

\documentclass[../../../include/open-logic-section]{subfiles}

\begin{document}

\olfileid[id]{fol}{com}{dls}
\olsection{Teorema L\"owenheim--Skolem}

Teorema L\"owenheim--Skolem menyatakan bahwa jika suatu teori mempunyai
model tak hingga, maka teori itu juga mempunyai model yang berhingga atau
!!{denumerable}. Salah satu konsekuensi langsung fakta ini ialah bahwa
logika orde pertama tidak dapat menyatakan bahwa ukuran !!a{structure}
bersifat !!{nonenumerable}: setiap !!{sentence} atau himpunan
!!{sentence} yang dipenuhi dalam semua !!{structure} !!{nonenumerable}
juga dipenuhi dalam suatu struktur !!{enumerable}.

\begin{thm}
\ollabel{thm:downward-ls} Jika $\Gamma$ konsisten, maka himpunan itu
mempunyai !!a{enumerable} model; yakni, dapat dipenuhi dalam
!!a{structure} yang domainnya berhingga atau !!{denumerable}.
\end{thm}

\begin{proof}
Jika $\Gamma$ konsisten, !!{structure}~$\Struct M$ yang dihasilkan oleh
pembuktian teorema kelengkapan mempunyai domain $\Domain{M}$ yang tidak
lebih besar daripada himpunan suku tertutup dari bahasa~$\Lang{L'}$.
Jadi, $\Struct M$ berhingga atau !!{denumerable}.
\end{proof}

\begin{thm}
\ollabel{noidentity-ls} Jika $\Gamma$ adalah himpunan konsisten dari
!!{sentence} dalam bahasa logika orde pertama tanpa identitas, maka
himpunan itu mempunyai !!a{denumerable} model; yakni, dapat dipenuhi
dalam !!a{structure} yang domainnya tak hingga dan !!{enumerable}.
\end{thm}

\begin{proof}
Jika $\Gamma$ konsisten dan tidak memuat kalimat yang mengandung
identitas, maka !!{structure}~$\Struct M$ yang dihasilkan oleh pembuktian
teorema kelengkapan mempunyai domain $\Domain{M}$ yang identik dengan
himpunan suku tertutup dari bahasa~$\Lang L'$. Jadi, $\Struct{M}$
!!{denumerable}: himpunan suku tertutup tersebut merupakan himpunan
bagian terhitung dari $\Trm[L']$ dan memuat tak hingga banyak konstanta
Henkin.
\end{proof}

\begin{ex}[Paradoks Skolem]
Teori himpunan Zermelo--Fraenkel~$\Log{ZFC}$ merupakan kerangka kerja yang
sangat kuat; di dalamnya, praktis semua pernyataan matematis dapat
dinyatakan, termasuk fakta-fakta mengenai ukuran himpunan. Sebagai
contoh, $\Log{ZFC}$ dapat membuktikan bahwa himpunan~$\Real$ dari bilangan
real bersifat !!{nonenumerable}; teori tersebut dapat membuktikan Teorema
Cantor bahwa himpunan kuasa sebarang himpunan lebih besar daripada
himpunan itu sendiri; dan sebagainya. Jika $\Log{ZFC}$ konsisten, semua
modelnya tak hingga. Lebih jauh, semua model tersebut memuat !!{element}
yang menurut teori bersifat !!{nonenumerable}; salah satu contohnya ialah
anggota yang keberadaannya membuat benar teorema~$\Log{ZFC}$ bahwa
himpunan kuasa bilangan asli ada. Berdasarkan Teorema
L\"owenheim--Skolem, $\Log{ZFC}$
juga mempunyai model-model !!{enumerable}---model yang memuat himpunan
``!!{nonenumerable}'' tetapi yang sendiri bersifat !!{enumerable}.
\end{ex}

\end{document}
